\section{The Four Faces Of The Number}

The backward walk now carries a bracketed number. Its lower face is charge,
the loop count. Its upper face is mass, the second-difference strain. But the
number is not exhausted by those two projections. The same loop also carries
a phase face, because orientation must be decided when it is transported, and
a sameness face, because the completed return touches the needle. The number
has four faces: charge, mass, phase, and sameness.

The first discipline is unity. These faces are not four substances placed in
one box. They are four readings of one carried loop. If the grammar allowed
charge to bring its own number, mass its own number, phase its own number,
and sameness its own number, the single-invariant work of the earlier
chapters would be lost. The loop would become a warehouse. The grammar
forbids that. It permits one number to be read by different faces because
the loop has returned with enough structure to support those readings.

Charge is the projection face nearest the count. It asks how many completed
traversals have returned with a signed orientation. Projection is cheap here
because the count has already been carried by loop closure. The charge face
does not decide a new orientation; it reads the tally that closure made
available. Its cost is the cost already paid by returning to the reference.

Mass is the other projection face. It asks how the loop strains when
variation passes null and threshold into response. This face is also a
projection, because the strain has already entered the bracket as the upper
reading. To read mass is to look at the carried second difference, not to
invent a second ledger. Projection is cheap only because the earlier section
paid the price of three trips: baseline, threshold, response.

Phase is more expensive. It is not merely a projection from the bracket. It
asks how the carried orientation is to be decided. A loop may return with an
orientation that can be read one way over the flat baseline and another way
over a tilted comparison. The phase face must therefore ask which decision
the current reading permits. It carries the two-fold sign possibility without
allowing either sign to become native before the baseline has been named.

This is why phase belongs between charge and sameness. Charge counts closure.
Phase reads how closure is oriented. A charge count without phase would know
that a loop returned but not how the return faced the reader. A phase without
charge would carry orientation without a completed traversal to bill. The
phase face is the decision cost of making orientation readable.

Sameness is the most delicate face. It is not a projection and not a free
decision. It is the identification face. It may touch the needle only because
the backward walk has returned to the sanctioned standing place after the
collapse of bare truth. The sameness face says that this completed return is
identified with the needle for this reading. It does not say that all
differences have disappeared. It says that the loop has closed through the
one bounded identity the grammar has earned.

The costs are therefore ordered. Charge and mass are free in the limited
sense that they are projections from the bracket already carried by the loop:
lower count and upper strain. Phase costs a decision, because the reader must
know which orientation the current baseline permits. Sameness costs the
quotient, because identity is never free in this book. To know a face is to
pay the cost appropriate to that face.

This cost accounting prevents two mistakes. The first mistake would treat
the four faces as synonyms. They are not. A loop count is not a second
difference. A phase decision is not an identity quotient. The second mistake
would treat the faces as separate entities. They are not that either. They
belong to one loop number, and each face is legitimate only because the loop
has carried the others far enough to keep the reading honest.

A simple ledger image helps. Suppose a completed route has returned to its
mark. One reader asks how many returns have occurred. That is charge. Another
asks how the route's response bends when perturbed past threshold. That is
mass. Another asks which way the route is facing when it returns through the
chosen comparison. That is phase. Another asks whether the returned mark can
stand as the same mark for the next act. That is sameness. The route is one;
the readings are four.

The phase face now becomes the next pressure point. Once orientation is
decision-bearing, the grammar must say how the sign split is read. The
imaginary language belongs here only if disciplined. The quarter-turn is not
a constructed object smuggled into the ledger. It is the reading by which the
phase face carries a bounce. Twice across that phase cost, the sign changes.
The next section follows that bounce without overclaiming it.
